% Chapitre 5 : l'assistant et son contexte.
\begin{figure}[tbp]
\centering
\begin{tikzpicture}[x=1cm,y=1cm,
  ctx/.style={boite, minimum width=3.1cm, text width=2.9cm}]
  \node[boite, minimum width=3.2cm, minimum height=1.3cm, draw=tonregle, fill=tonregle!7,
    line width=0.8pt, font=\sffamily\small] (s) at (0,0)
    {\textbf{Session}\\[-1pt]\footnotesize la conversation en cours};
  \node[ctx] (mcp) at (0,2.1) {\textbf{Serveurs MCP}\\[-1pt]\scriptsize issues, navigateur,\\Lean, articles};
  \node[ctx] (skills) at (-4.7,1.05) {\textbf{Skills}\\[-1pt]\scriptsize une procédure
    écrite par tâche, lue au besoin};
  \node[ctx] (mem) at (-4.7,-1.05) {\textbf{Mémoire}\\[-1pt]\scriptsize un fait par
    fichier, daté, relu au début};
  \node[ctx] (sub) at (4.7,0) {\textbf{Sous-agents}\\[-1pt]\scriptsize ne voient que leur
    consigne};
  \node[ctx] (repo) at (0,-2.2) {\textbf{Dépôt}\\[-1pt]\scriptsize fichiers, commandes,\\commits};
  \node[ctx] (cloud) at (4.7,-2.2) {\textbf{Session dans le cloud}\\[-1pt]\scriptsize
    clone le dépôt, pousse ses commits};
  \draw[fleche] (skills) -- (s);
  \draw[fleche, <->] (mem) -- (s);
  \draw[fleche, <->] (s) -- (mcp);
  \draw[fleche, <->] (s) -- (repo);
  \draw[fleche, <->] (cloud) -- (repo);
  \draw[fleche] ([yshift=2mm]s.east) -- node[etiq, above] {consigne} ([yshift=2mm]sub.west);
  \draw[fleche, tonregle] ([yshift=-2mm]sub.west) -- node[etiq, below, text=tonregle]
    {rapport,\\à vérifier} ([yshift=-2mm]s.east);
\end{tikzpicture}
\caption{L'assistant et son contexte. La session ne garde rien d'une fois sur l'autre :
ce qu'elle sait vient des skills, de la mémoire et du dépôt. Les sous-agents et les
sessions dans le cloud ne voient que leur consigne et le dépôt.}
\label{fig:assistant}
\end{figure}
